\documentclass{UAmathtalk}

\author{Frank Wagner}
\address{Cornell University}
\urladdr{https://sites.google.com/view/francis-wagner}
\title{Conjugator Lengths and Isoperimetric~Functions}
\date{Thursday, September 24, 2026}

\renewcommand*{\where}{Massry B008}
\renewcommand*{\abstractsize}{\normalsize}


\begin{document}

\maketitle

\begin{abstract}
The Dehn function of a finitely presented group is the smallest isoperimetric function of the Cayley complex associated to one of its finite presentations.  This invariant gives a geometric interpretation of a fundamental algorithmic notion, providing an upper bound on the complexity of the group’s Word Problem.

The conjugator length function of a finitely generated group is the function~$f$ so that $f(n)$ is the minimal upper bound on the length of a word realizing the conjugacy of two words of length at most~$n$.  Analogous to above, this invariant gives a geometric interpretation of an algorithmic notion, providing a measure for the complexity of a direct approach to the group’s Conjugacy Problem.

We show that these two notions are independent in a very strong sense, exhibiting finitely presented groups whose Dehn and conjugator length functions are equivalent to just about any `reasonable' pair of functions.  In particular, we find a finitely presented group with quadratic conjugator length function and undecidable Word Problem.  Moreover, we take a large step in describing all functions which can be realized as the conjugator length function of a finitely presented group.

This is joint work with Conan Gillis.
\end{abstract}

\end{document}
